% !PRINTED 27
\weberpage{27}
%<<CONT
\gesperrt{The divisor of a product of two whole functions is equal to
the product of the divisors of the two functions.}

For if $PA$ and $QB$ are two whole functions with divisors $P$ and
$Q$, then $A$ and $B$ are primitive functions. Hence $AB=C$ is also a
primitive function, and $PQ$ is the divisor of the function
$PA\cdot QB=PQC$.

Insofar as it refers to functions of one variable, we can give the
theorem the following form, without changing its content essentially;
in this form it is especially useful. Let
\[
 \varphi(x)=x^m+\alpha_1x^{m-1}+\alpha_2x^{m-2}+\cdots+\alpha_m,
\]
\[
 \psi(x)=x^n+\beta_1x^{n-1}+\beta_2x^{n-2}+\cdots+\beta_n
\]
\gesperrt{be two whole rational functions in which the highest powers
of $x$ have coefficient $1$, while the remaining coefficients are
rational numbers. Then, in the product}
\[
 \varphi(x)\psi(x)=x^{m+n}+\gamma_1x^{m+n-1}
 +\gamma_2x^{m+n-2}+\cdots+\gamma_{m+n}
\]
\gesperrt{the coefficients $\gamma$ cannot all be whole numbers if the
coefficients $\alpha,\beta$ in $\varphi(x)$ and $\psi(x)$ are not all
whole numbers.}

For let the smallest common denominator of the coefficients $\alpha$
of $\varphi$ be denoted by $a_0$, and that of the coefficients $\beta$
of $\psi$ by $b_0$. Then $a_0\varphi(x)=A(x)$ and
$b_0\psi(x)=B(x)$ are primitive functions of $x$. If the coefficients
$\gamma$ were whole numbers, the product $a_0b_0\varphi(x)\psi(x)$
would have divisor $a_0b_0$ and, if $a_0$ and $b_0$ were not both
equal to $1$, would therefore not be primitive; but this would
contradict the theorem proved above.

\begin{center}
{\large \S.\ 3.}

\vspace{0.5em}
{\large\gesperrt{Division.}}
\end{center}

Let, as before,
\begin{equation}
\begin{aligned}
 A&=A(x)=a_0x^m+a_1x^{m-1}+\cdots,\\
 B&=B(x)=b_0x^n+b_1x^{n-1}+\cdots
\end{aligned}
\tag{1}
\end{equation}
be two whole rational functions of $x$; but now suppose
%CONT>>
